\vspace{-5pt}
\section{Conclusion}
\vspace{-5pt}
In this work, we revisited personalized memory retrieval through the lens of dual-process theory. Existing methods are limited to one-shot retrieval based on similarity \emph{Familiarity}, whereas we introduce \emph{Recollection} as a deliberate stepwise retrieval mechanism. Building on this, we proposed \textbf{RF-Mem}, a familiarity uncertainty-driven framework that adaptively switches between them. Experiments show that RF-Mem achieves robust gains across both generation and retrieval tasks, and scales reliably to million-entry corpora, demonstrating the importance of integrating deliberate recollection into memory retrieval for personalizing LLM. 
% Looking forward, our findings open avenues for integrating cognitive-inspired retrieval strategies into broader retrieval-augmented LLM frameworks, enabling more adaptive and context-sensitive personalization at scale.
% In this work, we revisited personalized memory retrieval through the lens of dual-process theory. Existing approaches either rely solely on fast, similarity-based retrieval (\emph{Familiarity}) or adopt heavy-weight expansion mechanisms that resemble \emph{Recollection}, but neither accounts for the dynamic interplay between the two. We proposed \textbf{RF-Mem}, an entropy-driven framework that adaptively switches between Familiarity and Recollection according to retrieval uncertainty. Experiments on PersonaMem and PersonaBench demonstrate that RF-Mem achieves stable improvements across both factual and reasoning-intensive tasks, and remains robust under scaling to million-entry corpora. These results suggest that uncertainty-aware regulation is key to bridging efficiency and diagnostic depth in personalized memory systems. Looking forward, our findings open avenues for integrating cognitive-inspired switching strategies into broader retrieval-augmented LLM frameworks, enabling more adaptive and context-sensitive personalization at scale.